
\documentclass[a4paper,11pt]{article}
\usepackage[a4paper, margin=8em]{geometry}

% usa i pacchetti per la scrittura in italiano
\usepackage[french,italian]{babel}
\usepackage[T1]{fontenc}
\usepackage[utf8]{inputenc}
\frenchspacing 

% usa i pacchetti per la formattazione matematica
\usepackage{amsmath, amssymb, amsthm, amsfonts}

% usa altri pacchetti
\usepackage{gensymb}
\usepackage{hyperref}
\usepackage{standalone}

% imposta il titolo
\title{Appunti Elettronica Digitale}
\author{Luca Seggiani}
\date{2026}

% imposta lo stile
% usa helvetica
\usepackage[scaled]{helvet}
% usa palatino
\usepackage{palatino}
% usa un font monospazio guardabile
\usepackage{lmodern}

\renewcommand{\rmdefault}{ppl}
\renewcommand{\sfdefault}{phv}
\renewcommand{\ttdefault}{lmtt}

% disponi il titolo
\makeatletter
\renewcommand{\maketitle} {
	\begin{center} 
		\begin{minipage}[t]{.8\textwidth}
			\textsf{\huge\bfseries \@title} 
		\end{minipage}%
		\begin{minipage}[t]{.2\textwidth}
			\raggedleft \vspace{-1.65em}
			\textsf{\small \@author} \vfill
			\textsf{\small \@date}
		\end{minipage}
		\par
	\end{center}

	\thispagestyle{empty}
	\pagestyle{fancy}
}
\makeatother

% disponi teoremi
\usepackage{tcolorbox}
\newtcolorbox[auto counter, number within=section]{theorem}[2][]{%
	colback=blue!10, 
	colframe=blue!40!black, 
	sharp corners=northwest,
	fonttitle=\sffamily\bfseries, 
	title=~\thetcbcounter: #2, 
	#1
}

% disponi definizioni
\newtcolorbox[auto counter, number within=section]{definition}[2][]{%
	colback=red!10,
	colframe=red!40!black,
	sharp corners=northwest,
	fonttitle=\sffamily\bfseries,
	title=~\thetcbcounter: #2,
	#1
}

% U.D.M
\newcommand{\amp}{\ensuremath{\mathrm{A}}}
\newcommand{\volt}{\ensuremath{\mathrm{V}}}
\newcommand{\meter}{\ensuremath{\mathrm{m}}}
\newcommand{\second}{\ensuremath{\mathrm{s}}}
\newcommand{\farad}{\ensuremath{\mathrm{F}}}
\newcommand{\henry}{\ensuremath{\mathrm{H}}}
\newcommand{\siemens}{\ensuremath{\mathrm{S}}}

% circuiti
\usepackage{circuitikz}
\usetikzlibrary{babel}

% disegni
\usepackage{pgfplots}
\pgfplotsset{width=10cm,compat=1.9}

% disponi codice
\usepackage{listings}
\usepackage[table]{xcolor}

\definecolor{codegreen}{rgb}{0,0.6,0}
\definecolor{codegray}{rgb}{0.5,0.5,0.5}
\definecolor{codepurple}{rgb}{0.58,0,0.82}
\definecolor{backcolour}{rgb}{0.95,0.95,0.92}

\lstdefinestyle{codestyle}{
		backgroundcolor=\color{black!5}, 
		commentstyle=\color{codegreen},
		keywordstyle=\bfseries\color{magenta},
		numberstyle=\sffamily\tiny\color{black!60},
		stringstyle=\color{green!50!black},
		basicstyle=\ttfamily\footnotesize,
		breakatwhitespace=false,         
		breaklines=true,                 
		captionpos=b,                    
		keepspaces=true,                 
		numbers=left,                    
		numbersep=5pt,                  
		showspaces=false,                
		showstringspaces=false,
		showtabs=false,                  
		tabsize=2
}

\lstdefinestyle{shellstyle}{
		backgroundcolor=\color{black!5}, 
		basicstyle=\ttfamily\footnotesize\color{black}, 
		commentstyle=\color{black}, 
		keywordstyle=\color{black},
		numberstyle=\color{black!5},
		stringstyle=\color{black}, 
		showspaces=false,
		showstringspaces=false, 
		showtabs=false, 
		tabsize=2, 
		numbers=none, 
		breaklines=true
}

\lstdefinelanguage{javascript}{
	keywords={typeof, new, true, false, catch, function, return, null, catch, switch, var, if, in, while, do, else, case, break},
	keywordstyle=\color{blue}\bfseries,
	ndkeywords={class, export, boolean, throw, implements, import, this},
	ndkeywordstyle=\color{darkgray}\bfseries,
	identifierstyle=\color{black},
	sensitive=false,
	comment=[l]{//},
	morecomment=[s]{/*}{*/},
	commentstyle=\color{purple}\ttfamily,
	stringstyle=\color{red}\ttfamily,
	morestring=[b]',
	morestring=[b]"
}

% disponi sezioni
\usepackage{titlesec}

\titleformat{\section}
	{\sffamily\Large\bfseries} 
	{\thesection}{1em}{} 
\titleformat{\subsection}
	{\sffamily\large\bfseries}   
	{\thesubsection}{1em}{} 
\titleformat{\subsubsection}
	{\sffamily\normalsize\bfseries} 
	{\thesubsubsection}{1em}{}

% disponi alberi
\usepackage{forest}

\forestset{
	rectstyle/.style={
		for tree={rectangle,draw,font=\large\sffamily}
	},
	roundstyle/.style={
		for tree={circle,draw,font=\large}
	}
}

% disponi algoritmi
\usepackage{algorithm}
\usepackage{algorithmic}
\makeatletter
\renewcommand{\ALG@name}{Algoritmo}
\makeatother

% disponi numeri di pagina
\usepackage{fancyhdr}
\fancyhf{} 
\fancyfoot[L]{\sffamily{\thepage}}

\makeatletter
\fancyhead[L]{\raisebox{1ex}[0pt][0pt]{\sffamily{\@title \ \@date}}} 
\fancyhead[R]{\raisebox{1ex}[0pt][0pt]{\sffamily{\@author}}}
\makeatother

\begin{document}
% sezione (data)
\section{Lezione del 12-05-26}

% stili pagina
\thispagestyle{empty}
\pagestyle{fancy}

% testo
\subsection{Logica a pass transistor}
Vediamo adesso una categoria di famiglie logiche spesso associate alle porte MOS complementari, in particolare per realizzare funzioni come i \textit{multiplexer}. Alla base della logica pass transistor è l'uso degli interruttori (cioè i transistor) come elementi di passo.

Vediamo ad esempio la seguente configurazion:
\begin{center}
\begin{tikzpicture}[transform shape]
	% Paths, nodes and wires:
	\draw (-4, 1) -- (-2, 1);
	\draw (-4, -1) -- (-2, -1);
	\draw (-4, 1) to[open, v={$V_A=A$}, voltage=american] (-4, -1);
	\draw (-2, 1) to[switch, l={$B$}] (-1, 1);
	\draw (-0.995, 1) to[switch, l={$C$}] (0.005, 1);
	\draw (2, 1) to[american resistor, l={$R$}] (2, -1);
	\draw (2, -1) |- (-2, -1);
	\draw (0.005, 1) -- (2, 1);
	\draw (2, -1) -- (4, -1);
	\draw (2, 1) -- (4, 1);
	\draw (4, 1) to[open, v={$V_y=Y$}, voltage=american] (4, -1);
\end{tikzpicture}
\end{center}

In questo caso, i due interruttori ($B$ e $C$) in serie possono essere considerate come i due valori di una porta AND, con la $A$ che fa da terzo. Allo stesso modo, se avessimo messo $B$ e $C$ in parallelo, avremmo ottenuto una porta OR (con la $A$ sempre ad AND in quanto pilota l'intero circuito). 

La cosa importante in questi tipi di circuiti è di assicurare che il nodo di uscita $y$ non sia lasciato flottante, cioè non sia lasciato senza percorsi a ground o a $V_{DD}$.

\par\smallskip

Il problema a questo punto sarà come implementare l'interruttore.
\begin{itemize}
\item L'opzione più semplice è quella di usare un \textbf{NMOS}. Per verificare che un NMOS si comporti effettivamente come interruttore, inseriamolo in un circuito dove questo alimenta una capacità $C$.

\par\smallskip
\textit{(Carica con NMOS)}
\begin{enumerate}
	\item 
Poniamo la carica di $C$ come nulla a $t = 0$, e l'alimentazione come $V_{DD}$. Ora, in un circuito integrato (dove il drain non viene fissato a body), drain e source di un MOSFET sono completamente interscambiabili, per cui non ci interessa distinguere fra i due.

Applicando una tensione $V_{DD}$ al gate del MOSFET, avremo che la $V_{GS}$ si potrà calcolare come:
$$
V_{GS} = V_G - V_S = V_{DD} - V_{C} \implies V_{GS} = V_{DD} > V_T
$$
visto che il source è il nodo collegato al condensatore (è scarico, la corrente passa da lui a ground). All'inizio del ciclo di carica, $V_C = 0$ per cui $V_{DD} > V_T$ ed il transistor conduce. La $V_{DS}$ sarà ugualmente:
$$
V_{DS} = V_D - V_S = V_{DD} \geq V_{DD} - V_T \ \text{(saturo)}
$$
cioè saremo nella regione di saturazione del NMOS. Il risultato è che il condensatore viene caricato come se fosse collegato ad una sorgente di tensione $V_{DD}$;
	\item 
Questa situazione non può chiaramente continuare per sempre. Infatti, l'NMOS \textit{conduce} se viene rispettata la condizione:
$$
V_{GS} \geq V_T \ \text{(conduzione)}
$$
mentre è \textit{saturo} se viene rispettata la condizione:
$$
V_{DS} \geq V_{GS} - V_T \ \text{(saturazione)}
$$

Per la prima condizione (di conduzione) si ha che:
$$
V_G - V_S \geq V_T, \quad V_{DD} - V_C \geq V_T \implies V_C \leq V_{DD} - V_T
$$
Per la seconda (di saturazione), si ha invece che:
$$
V_D - V_S \geq V_G - V_S - V_T, \quad V_D \geq V_G - V_T, \quad V_{DD} \geq V_{DD} - V_T 
$$
cioè l'NMOS è sempre in saturazione. 

La corrente $i_D$ sarà quindi data da:
$$
i_D = K \left( V_{DS} - V_T \right)^2
$$
in quanto $V_{DS} = V_{GS}$, dove:
$$
V_{DS} = V_D - V_S = V_{DD} - V_C \geq V_T
$$

\end{enumerate}

Abbiamo quindi che il MOSFET si comporta effettivamente da interruttore, seppur non ideale (conduce per $V_C = 0$, ma non riesce a portare la $V_C$ fino a $V_{DD}$, cioè lascia un offset di $V_T$ oltre la quale non riesce più a condurre).
Risulta quindi asimmetrico in quanto riesce a chiudere perfettamente per scaricare a $0 \, \mathrm{V}$, ma non per caricare a $V_{DD}$. 

\par\smallskip
\textit{(Scarica con NMOS)}
\begin{enumerate}
	\item 
Possiamo confermare questa asimmetria vedendo situazione duale, dove la $V_C = V_{DD}$ (la capacità è gia carica) e vogliamo usare il MOSFET per scaricarla a ground. 
In questo caso, posta sempre $V_G = V_{DD}$, si ha:
$$
V_{GS} = V_G - V_S = V_{DD} - 0 \geq V_T 
$$
in quanto $V_S$ è sempre nulla (stiamo conducendo la scarica del condensatore verso ground), e il MOSFET è sempre in conduzione.

Per quanto riguarda la saturazione, abbiamo che la $V_{DS}$ è:
$$
V_{DS} = V_D - V_S = V_C > V_T
$$
e quindi siamo in saturazione per $V_C > V_T$, regime ohmico altrimenti.

Riassumiamo quanto detto: siamo sempre in conduzione, e la $V_{GS}$ non cambia in quanto la $V_S$ è fissa a 0. Questo significherà che seguiamo sempre la stessa caratteristica fino a $V_{DS} = 0 \, \mathrm{V}$, cioè fino alla scarica completa della capacità. 

\end{enumerate}

\item Vediamo cosa succede se si usa un \textbf{PMOS} nella stessa situazione. Anche qui, verifichiamo inserendo il PMOS nel circuito di carica di una capacità.

\par\smallskip
\textit{(Carica con PMOS)}
\begin{enumerate}
	\item
Poniamo la carica di $C$ come nulla a $t = 0$, e l'alimentazione a $V_{DD}$, come nel caso precedente.

Perchè il MOSFET conduca bisognamo di $V_G$ a $0 \, \mathrm{V}$, in quanto la $V_S$ (sul PMOS dalla parte dell'alimentazione) sarà $V_{DD}$, e quindi:
$$
V_{GS} = V_G - V_S = 0 - V_{DD} = -V_{DD} \leq V_T
$$
e siamo sempre in conduzione. Il fatto che siamo sempre in conduzione a $V_{GS}$ costante ci mette nella situazione in cui seguiamo un unica caratteristica, finche $V_{DS}$ non è nulla (quando $V_C$ è carico). 
\end{enumerate}

Il PMOS è quindi ideale nella situazione opposta all'NMOS: riesce a traspostare l'\textit{"uno pieno"}, cioè a caricare completamente fino a $V_{DD}$. Intuitivamente, il PMOS sarà duale all'NMOS anche nella scarica, e quindi si comporterà in maniera meno ideale:

\par\smallskip
\textit{(Carica con PMOS)}
\begin{enumerate}
	\item 
Iniziamo con la $V_C = V_{DD}$ e il PMOS collegato a massa al drain (chiamiamo il nodo drain tale solo perché è a massa, in quanto come l'NMOS il PMOS integrato è simmetrico). Segue che la $V_{GS}$ è:
$$
V_{GS} = V_G - V_S = 0 - V_C
$$
Quindi la $V_{GS}$ non è più costante nel tempo, ma anzi passa da $V_{DD}$ (carica iniziale della $C$) a $0 \, \mathrm{V}$. 

Perché ci sia conduzione sul transistor, ricordiamo che serve:
$$
V_{GS} \leq V_T \implies V_C \geq -V_T
$$
e riguardo alla saturazione:
$$
V_{DS} = V_D - V_S = V_{GS} = V_G - V_S 
$$
cioè il PMOS è sempre in saturazione.

	\item
Questa situazione non può chiaramente continuare all'infinito, per cui arriveremo ad un punto dove:
$$
V_{GS} = 0 - V_C > V_T \implies V_C < -V_T
$$
In questo punto il transistor non è più in conduzione, e la capacità non scaricherà oltre. Questo significherà che abbiamo trasmesso gli $0 \, \mathrm{V}$ fino a $- V_T$ e non oltre, trovandoci nella stessa situazione (ma duale) dell'NMOS in carica.
\end{enumerate}

\end{itemize}

Riassumendo, abbiamo ottenuto che:
\begin{itemize}
	\item L'NMOS è un ottimo interruttore per la transizione da $V_{DD}$ $\rightarrow$ $0 \, \mathrm{V}$;
	\item Il PMOS è un ottimo interruttore per la transizione da $0 \, \mathrm{V}$ $\rightarrow$ $V_{DD}$ ;
\end{itemize}

La soluzione per realizzare un interruttore ideale attraverso i transistori MOSFET sarà allora quella di usarne due, uno di tipo NMOS e uno di tipo PMOS, in parallelo. L'NMOS sarà pilotato dalla variabile di controllo dello switch (diciamo $B$), mentre il PMOS sarà pilotato dalla $\overline{B}$: questo perchè ricordiamo l'NMOS risponde a $V_G > V_T$ e il PMOS viceversa.

\begin{center}
\begin{tikzpicture}[transform shape]
	% Paths, nodes and wires:
	\node[pmos, rotate=90] at (0, -0.5){};
	\node[nmos, rotate=-90] at (0, 0.5){};
	\draw (-0.77, -0.5) -- (-0.77, 0.5);
	\draw (0.77, 0.5) |- (0.75, -0.5);
	\draw (-2, -0) -| (-0.79, -0.003);
	\draw (0.79, 0.003) -| (2, -0);
	\node[shape=rectangle, minimum width=0.465cm, minimum height=0.465cm] at (0, 1.75){} node[anchor=south, align=center, text width=0.077cm, inner sep=6pt] at (0, 1.5){B};
	\node[shape=rectangle, minimum width=0.465cm, minimum height=0.465cm] at (0, -1.75){} node[anchor=north, align=center, text width=0.077cm, inner sep=6pt] at (0, -1.5){$\overline{\text{B}}$};
\end{tikzpicture}
\end{center}

Questo componente viene anche detto \textit{pass gate}.

\subsubsection{Multiplexer 2:1 a pass gate}
Il circuito di esempio che portiamo per un multiplexer è il seguente:
\begin{center}
\begin{tikzpicture}[transform shape]
	% Paths, nodes and wires:
	\draw (-1, 1) to[switch, l={$C$}] (1, 1);
	\draw (-0.995, -1) to[switch, l={$\overline{\text{C}}$}] (1.005, -1);
	\draw (1, 1) -| (1, -0) -- (2, -0);
	\draw (1.005, -1) |- (1, -0);
	\node[shape=rectangle, minimum width=0.465cm, minimum height=0.465cm] at (-1.25, 1){} node[anchor=center, align=center, text width=0.077cm, inner sep=6pt] at (-1.25, 1){A};
	\node[shape=rectangle, minimum width=0.465cm, minimum height=0.465cm] at (-1.25, -1){} node[anchor=center, align=center, text width=0.077cm, inner sep=6pt] at (-1.25, -1){B};
	\node[shape=rectangle, minimum width=0.465cm, minimum height=0.465cm] at (2.25, -0){} node[anchor=center, align=center, text width=0.077cm, inner sep=6pt] at (2.25, -0){$y$};
\end{tikzpicture}
\end{center}

Vediamo chiaramente come pilotare i due interruttori, uno con $C$ e l'altro con $\overline{C}$, porta al comportamento di "isolamento" di una delle due entrate dei multiplexer.

Implementando ogni interruttore come un pass gate, otteniamo che il circuito si realizza con $4$ transistori, più $2$ per l'inverter della $C$, per cui $6$. Di contro, scegliendo di usare la tecnologia complementare, avremmo avuto da implementare la funzione:
$$
y = AC + B \overline{C}
$$
Questa è una funzione realizzata con due AND e una OR, per cui $4 \cdot 2 = 8$ MOS solo le per la circuiteria logica (sia PMOS che CMOS). Inoltre, abbiamo tutte le porte sia negate che non, per cui abbiamo bisogno di $3$ inverter, quindi $14$ transistor in totale. L'implementazione a pass gate è molto più conservativa in termini di transistori usati! 

\subsubsection{XOR a 2 ingressi a pass gate}
Vediamo come si può realizzare anche una XOR in maniera simile. Questa era infatti data dalla formula logica:
$$
y = A \overline{B} + \overline{A} B
$$
che si implementa in pass gate in maniera esattamente uguale alla precedente.

In questo caso usiamo $4$ transistor per i pass gate, più $2$ inverter, da cui $4 + 2 \cdot 2 = 8$ transistor. Come avevamo visto in 25.0.1, la stessa cosa fatta in tecnologia complementare usava $12$ transistor.


\end{document}
